
          %%%%% ~~~~~~~~~~~~~~~~~~~~ %%%%%

\chapter{Time-dependent perturbation theory and the interaction picture}
\setcounter{theorem}{0}
\setcounter{equation}{0}

%\cite{vanDerVaart1996}
%\cite{Kosorok2008}

%\renewcommand{\theenumi}{\alph{enumi}}
%\renewcommand{\labelenumi}{\textnormal{(\theenumi)}$\;\;$}
\renewcommand{\theenumi}{\roman{enumi}}
\renewcommand{\labelenumi}{\textnormal{(\theenumi)}$\;\;$}

          %%%%% ~~~~~~~~~~~~~~~~~~~~ %%%%%

We want to determine the time-evolution of a system under the time-varying Hamiltonian
\begin{equation*}
H(t) \;\; = \;\; H_{0} \,+\, H_{1}(t)
\end{equation*}

\vskip 0.5cm
\noindent
\textbf{Comments}
\begin{itemize}
\item
	We seek solution in the form of a one-parameter family of unitary operators
	\,$U(t)$\, which satisfies
	\begin{equation*}
	\dfrac{\d}{\d t}\,U(t) \;\; = \;\; -\,\i\,H(t)\,U(t)\,,
	\quad
	\textnormal{where \,$\i \,:=\, \sqrt{-1}$}
	\end{equation*}
\item
	We assume we know that the time-evolution of a system under the constant (with respect to time)
	Hamiltonian $H_{0}$ is given by
	\begin{equation*}
	U_{0}(t) \;\; = \;\; \exp\left(\,-\,\i\,t\,H_{0}\,\right),
	\end{equation*}
	which is determined as the solution of:
	\begin{equation*}
	\left\{\;\;
		\begin{array}{ccc}
		\dfrac{\d}{\d t}\,U_{0}(t) & = & -\,\i\,H_{0}\,U_{0}(t)
		\\
		\overset{{\color{white}1}}{U_{0}(0)} & = & \textnormal{identity operator}
		\end{array}
		\right.
	\end{equation*}
\end{itemize}

          %%%%% ~~~~~~~~~~~~~~~~~~~~ %%%%%

\vskip 0.5cm
\section{The interaction picture}

          %%%%% ~~~~~~~~~~~~~~~~~~~~ %%%%%

The situation where we time-evolve the operator \,$A$\, according to
\begin{equation*}
A(t) \;\; = \;\; U_{0}(-t) \cdot A \cdot U_{0}(t)
\end{equation*}
and the state \,$x(t)$\, according to
\begin{equation*}
x(t) \;\; = \;\; U_{0}(-t) \cdot U(t) \cdot x
\end{equation*}
is called the \textbf{interaction picture}.
Note that the interaction picture is consistent with both the Schr\"odinger and the Heisenberg pictures:
\begin{eqnarray*}
\left\langle\; \overset{{\color{white}.}}{x(t)} \,,\, A(t) \cdot x(t) \;\right\rangle
& = &
	\left\langle\;
		\overset{{\color{white}.}}{U_{0}(-t)} \cdot U(t) \cdot x
		\;\,,\;
		U_{0}(-t) \cdot A \cdot U_{0}(t) \cdot U_{0}(-t) \cdot U(t) \cdot x
		\;\right\rangle
\\
& = &
	\left\langle\;
		\overset{{\color{white}.}}{U(t)} \cdot x
		\;\,,\;
		A \cdot U(t) \cdot x
		\;\right\rangle
	\quad\quad
	\textnormal{(Schr\"odinger picture)}
\\
& = &
	\left\langle\;
		\overset{{\color{white}.}}{U(-t)} \cdot U(t) \cdot x
		\;\,,\;
		U(-t) \cdot A \cdot U(t) \cdot x
		\;\right\rangle
\\
& = &
	\left\langle\;
		x
		\;\,,\;
		\overset{{\color{white}.}}{U(-t)} \cdot A \cdot U(t) \cdot x
		\;\right\rangle
	\quad\quad
	\textnormal{(Heisenberg picture)}
\end{eqnarray*}

          %%%%% ~~~~~~~~~~~~~~~~~~~~ %%%%%

\vskip 0.5cm
\section{The Dyson expansion}

          %%%%% ~~~~~~~~~~~~~~~~~~~~ %%%%%

\begin{proposition}[Dyson expansion]
\mbox{}\vskip 0.2cm
\noindent
Define \,$V(t) \, := \, U_{0}(t)^{-1} \cdot U(t) \, = \, U_{0}(-t) \cdot U(t)$.\,
Then, the following statements are true:
\begin{enumerate}
\item
	We have:\quad $U(t) \,= \, U_{0}(t) \cdot V(t)$
\item
	$V(t)$\, satisfies:
	\begin{equation*}
	\dfrac{\d}{\d t}\,V(t) \;\; = \;\; -\,\i\,\widetilde{H}_{1}(t)\,V(t)\,,
	\end{equation*}	
	where
	\begin{equation*}
	\widetilde{H}_{1}(t) \;\; := \;\; U_{0}(-t) \cdot H_{1}(t) \cdot U_{0}(t)
	\end{equation*}	
\item
	$V(t)$ can be expressed as:
	\begin{equation*}
	V(t) \;\; = \;\; \mathbf{1} \, + \, \overset{\infty}{\underset{n = 1}{\sum}}\,W_{n}(t)\,,
	\end{equation*}
	where
	\begin{equation*}
	W_{n}(t) \;\; = \;\; (-\i)^{n}\cdot
		\int_{0}^{t}\d\theta_{1} \, \int_{0}^{\theta_{1}}\d\theta_{2} \;\, \cdots \,  \int_{0}^{\theta_{n-1}}\d\theta_{n}
			\left(\,
				\overset{{\color{white}1}}{\widetilde{H}_{1}(\theta_{1})}
				\cdot
				\widetilde{H}_{1}(\theta_{2})
				\,\cdots\,
				\widetilde{H}_{1}(\theta_{n})
				\,\right),
	\;\;
	\end{equation*}
	for \,$n = 1, 2, \ldots$
\item
	For \,$n = 1, 2, \ldots$\,, \,$W_{n}(t)$\, can be expressed as:
	\begin{equation*}
	W_{n}(t) \;\; = \;\; \dfrac{(-\i)^{n}}{n!}\cdot
		\int_{0}^{t}\d\theta_{1} \, \int_{0}^{t}\d\theta_{2} \;\, \cdots \,  \int_{0}^{t}\d\theta_{n}\;
			\mathcal{T}\!\left(\,
				\overset{{\color{white}1}}{\widetilde{H}_{1}(\theta_{1})}
				\cdot
				\widetilde{H}_{1}(\theta_{2})
				\,\cdots\,
				\widetilde{H}_{1}(\theta_{n})
				\,\right),
	\end{equation*}
	where \,$\mathcal{T}(\,\cdot\,)$ is the \textbf{chronological operator}, defined as follows:
	\begin{equation*}
	\mathcal{T}\!\left(\,
		\overset{{\color{white}1}}{\widetilde{H}_{1}(\theta_{1})}
		\cdot
		\widetilde{H}_{1}(\theta_{2})
		\,\cdots\,
		\widetilde{H}_{1}(\theta_{n})
		\,\right)
	\;\; := \;\;
		\overset{{\color{white}1}}{\widetilde{H}_{1}(\theta_{\sigma(1)})}
		\cdot
		\widetilde{H}_{1}(\theta_{\sigma(2)})
		\,\cdots\,
		\widetilde{H}_{1}(\theta_{\sigma(n)})\,,
	\end{equation*}
	where \,$\sigma \in S_{n}$\, is any permutation of $\{1,2,\cdots,n\}$ such that
	\begin{equation*}
	\theta_{\sigma(1)} \; \leq \; \theta_{\sigma(2)} \; \leq \; \cdots \; \leq \; \theta_{\sigma(n)}
	\end{equation*}
\item
	For \,$H_{1}(t) \,=\, \kappa\,H_{I}$,\,
	\begin{equation*}
	V(t) \;\; = \;\; \mathbf{1} \, + \,
		\overset{\infty}{\underset{n = 1}{\sum}}\, \dfrac{(-\i\,\kappa)^{n}}{n!}
		\cdot
		\int_{0}^{t}\d\theta_{1} \, \int_{0}^{t}\d\theta_{2} \;\, \cdots \,  \int_{0}^{t}\d\theta_{n}\;
			\mathcal{T}\!\left(\,
				\overset{{\color{white}1}}{H_{I}(\theta_{1})}
				\cdot
				H_{I}(\theta_{2})
				\,\cdots\,
				H_{I}(\theta_{n})
				\,\right),
	\end{equation*}
	where \,$H_{I}(t) \, := \, U_{0}(-t) \cdot H_{I} \cdot U_{0}(t)$.
\end{enumerate}
\end{proposition}

          %%%%% ~~~~~~~~~~~~~~~~~~~~ %%%%%

\vskip 0.5cm
\section{The Lehmann-Symanzik-Zimmermann reduction formula}

          %%%%% ~~~~~~~~~~~~~~~~~~~~ %%%%%

\begin{eqnarray*}
&&
	\left(\,
	    \overset{m}{\underset{i=1}{\prod}}\;
	    \dfrac{\i\sqrt{Z}}{k_{i}^{2}-m^{2}}
	    \right)
	\!
	\left(\,
	    \overset{n}{\underset{j=1}{\prod}}\;
	    \dfrac{\i\sqrt{Z}}{p_{j}^{2}-m^{2}}
	    \right)
	\!
	{\color{red}
	\left\langle\,
	    \mathbf{p}_{1} \cdots \mathbf{p}_{n}
	    \,\vert\,
	    \i\,T
	    \,\vert\,
	    \mathbf{k}_{1} \cdots \mathbf{k}_{m}
	    \,\right\rangle
	}
\\
& = &
    \bigints\;
	    \overset{m}{\underset{i=1}{\prod}}\;e^{-\i\,k_{i} \cdot x_{i}}
	    \;
	    \overset{n}{\underset{j=1}{\prod}}\;e^{+\i\,p_{j} \cdot y_{j}}
	    \,
    	{\color{red}
		\left\langle\,
		    0
		    \,\left\vert\,
		    \mathbb{T}\left\{\,
		        \overset{{\color{white}.}}{\phi(x_{1})\cdots\phi(x_{1})}
		        \,
		        \phi(y_{1})\cdots\phi(y_{1})
		        \,\right\}
		    \,\right\vert\,
		    0
		    \,\right\rangle
		}
		\;
		\d^{4}x_{1} \cdots \d^{4}x_{m}
		\;
		\d^{4}y_{1} \cdots \d^{4}y_{n}
\end{eqnarray*}

          %%%%% ~~~~~~~~~~~~~~~~~~~~ %%%%%

\vskip 0.5cm
\noindent
\begin{itemize}
\item
    Hamiltonian:\; $H \,=\, H_{0} + H_{I}$,\, where
    \,$H_{0}$\, is a free-field Hamiltonian, and
    \,$H_{I}$\, is a small perturbation that describes interaction.
\item
    $U_{0}(t) \, := \, e^{-\,\i tH_{0}}$,\,
    \quad
    $U(t) \, := \, e^{-\,\i tH}$
\item
    Define:
    \,$
    \phi_{I}(t,\mathbf{x})
    \, := \,
        U_{0}(-t) \cdot U(t) \cdot \phi(t,\mathbf{x}) \cdot U(-t) \cdot U_{0}(t)
    \, = \,
        e^{+\i\,tH_{0}} \cdot e^{-\i\,tH} \cdot \phi(t,\mathbf{x}) \cdot e^{+\i\,tH} \cdot e^{-\i\,tH_{0}}
    $.\,
    \vskip 0.05cm
    Then, in particular, we have:
    \,$
    \phi_{I}(t,\mathbf{x})
    \, = \,
        e^{+\i\,tH_{0}} \cdot e^{-\i\,tH} \cdot \phi(t,\mathbf{x}) \cdot e^{+\i\,tH} \cdot e^{-\i\,tH_{0}}
    $\,
\end{itemize}

          %%%%% ~~~~~~~~~~~~~~~~~~~~ %%%%%

\vskip 0.5cm
\section{Cross sections \& differential cross sections of particle scattering}

          %%%%% ~~~~~~~~~~~~~~~~~~~~ %%%%%

\begin{lemma}
\mbox{}
\vskip 0.01cm
\noindent
Suppose:
\begin{itemize}
\item
    Two cylindrial packets of protons are made to collide head on
    (travelling in exactly opposite directions in a particle collider).
\item
    $A$\, is the common geometric cross-sectional area each of the packets,
    and
    \,$L$\, is their common length.
    The common volume of the two packets is thus \,$V = A\,L$.\,
\item
    $N_{1}$\, and \,$N_{2}$\, are the number of protons 
    in the first and second packets of protons, respectively.
\item
    The number density \,$n_{1}$\, of protons in the first packet is constant in space and time.
    Similarly,
    the number density \,$n_{2}$\, of protons in the second packet is constant in space and time.
\item
    $v_{1}$\, and \,$v_{2}$\, are the speeds (as observed in the reference frame of the laboratory)
    of the first and second packets of protons, respectively.
    Hence, the relative speed of the two packets of protons is
    \,$\vert\,v_{1} - v_{2}\,\vert$.\,
\end{itemize}
Let
\,$V_{\textnormal{overlap}}(t) \subset \Re^{3}$\,
denote the overlap at time \,$t \in \Re$\, of the two packets.
\vskip 0.2cm
\noindent
\!We assume:
There exists a
{\color{red}(proportionality) constant \,$\sigma > 0$\,}
such that the
\textbf{expected number \,$r(t,x)$\, at time \,$t \in \Re$\, and at \,$x \in \Re^{3}$
of interactions per unit time and per unit volume}\,
is given by:
\begin{equation*}
{\color{red}r(t,x)}
\;\; = \;\;
    \left\{\begin{array}{cl}
        {\color{red}\sigma \cdot \vert\,v_{1} - v_{2}\,\vert \cdot n_{1} n_{2}}\,,
        &
        \textnormal{for \,$t \in \Re$\, such that
            \,$V_{\textnormal{overlap}}(t) \,\neq\, \varemptyset$\,
            and \,$x \in V_{\textnormal{overlap}}(t)$\,
            }
        \\
        \overset{{\color{white}1}}{0}\,,
        &
        \textnormal{otherwise}
        \end{array}\right.
\end{equation*}
(Note that the quantity
\,$\sigma \cdot \vert\,v_{1} - v_{2}\,\vert \cdot n_{1} n_{2}$\,
above is in fact a constant, independent of \,$t$\, and \,$x$.)
\vskip 0.2cm
\noindent
Then, the following equality is true:
\begin{equation*}
\left(\!\begin{array}{c}
    \textnormal{total number of}
    \\
    \textnormal{interaction events}
    \end{array}\!\right)
\;\; = \;\;
    \sigma \cdot \dfrac{ N_{1} N_{2} }{A}
\end{equation*}
\end{lemma}
\proof
We choose the reference time
\,$t = 0$\,
to be the exact moment when the centres of the two cylindrical packets coincide
(equivalently, when the two packets perfecrtly overlap each other).
We will use the slight abuse of notation to denote the volume
of the time-dependent overlap
\,$V_{\textnormal{overlap}}(t)$\,
also as
\,$V_{\textnormal{overlap}}(t)$.\,

\vskip 0.3cm
\noindent
Let
\,$R(t)$\,
be the total number of interactions per unit time (over all of space).
Then, we have:
\begin{eqnarray*}
R(t)
& = &
    \int_{\Re^{3}}\;
        r(t,x)
        \;\d^{3}x
\;\; = \;\;
    \int_{V_{\textnormal{overlap}}(t)}\;
        \sigma \cdot \vert\,v_{1} - v_{2}\,\vert \cdot n_{1} n_{2}
        \;\d^{3}x
\\
& = &
    \sigma \cdot \vert\,v_{1} - v_{2}\,\vert \cdot n_{1} n_{2} \cdot V_{\textnormal{overlap}}(t)
\end{eqnarray*}
Next, observe that
\begin{equation*}
V_{\textnormal{overlap}}(t)
\;\; = \;\;
    \left\{\begin{array}{cl}
        A\left(L \,\overset{{\color{white}.}}{-}\, \vert\,v_{1}-v_{2}\,\vert\cdot\vert\,t\,\vert\right)\,, & \textnormal{for \;$t \,\in\, \left[\,-\,\dfrac{L}{\vert\,v_{1}-v_{2}\,\vert}\,,\,\dfrac{L}{\vert\,v_{1}-v_{2}\,\vert}\,\right]$}
        \\
        0\,, & \textnormal{otherwise}
        \end{array}\right.
\end{equation*}
Hence,
\begin{eqnarray*}
\left(\!\begin{array}{c}
    \textnormal{number of}
    \\
    \textnormal{interaction events}
    \end{array}\!\right)
& = &
    \bigintsss_{-\infty}^{+\infty}\;
        R(t)
        \;\d t
\;\; = \;\;
    \bigintsss_{-\infty}^{+\infty}\;
        \sigma \cdot \vert\,v_{1} - v_{2}\,\vert \cdot n_{1} n_{2} \cdot V_{\textnormal{overlap}}(t)
        \;\d t
\\
& = &
    \sigma \cdot \vert\,v_{1} - v_{2}\,\vert \cdot n_{1} n_{2} \cdot
    \bigintsss_{-L/\vert\,v_{1}-v_{2}\,\vert}^{+L/\vert\,v_{1}-v_{2}\,\vert}\;
        V_{\textnormal{overlap}}(t)
        \;\d t
    % \bigintsss_{-L/\vert\,v_{1}-v_{2}\,\vert}^{+L/\vert\,v_{1}-v_{2}\,\vert}\;
    %     A (L \,-\, \vert\,v_{1}-v_{2}\,\vert\cdot\vert\,t\,\vert)
    %     \;\d t
\\
& = &
    \sigma \cdot \vert\,v_{1} - v_{2}\,\vert \cdot n_{1} n_{2} \cdot
    \bigintsss_{-L/\vert\,v_{1}-v_{2}\,\vert}^{+L/\vert\,v_{1}-v_{2}\,\vert}\;
         A\left(L \,\overset{{\color{white}.}}{-}\, \vert\,v_{1}-v_{2}\,\vert\cdot\vert\,t\,\vert\right)
        \;\d t
\\
& = &
    \sigma \cdot \vert\,v_{1} - v_{2}\,\vert \cdot n_{1} n_{2} \cdot
    \left(\,
	    AL \cdot \dfrac{2\,L}{\vert\,v_{1}-v_{2}\,\vert}
	    \; - \;
	    A \cdot \vert\,v_{1}-v_{2}\,\vert
	    \bigintsss_{-L/\vert\,v_{1}-v_{2}\,\vert}^{+L/\vert\,v_{1}-v_{2}\,\vert}\;
	        \vert\,t\,\vert
	        \;\d t
	    \,\right)
\\
& = &
    \sigma \cdot \vert\,v_{1} - v_{2}\,\vert \cdot n_{1} n_{2} \cdot
    \left(\,
	    2\cdot\dfrac{AL^{2}}{\vert\,v_{1}-v_{2}\,\vert}
	    \; - \;
	    2\,A \cdot \vert\,v_{1}-v_{2}\,\vert
	    \bigintsss_{0}^{+L/\vert\,v_{1}-v_{2}\,\vert}\;
	        t
	        \;\d t
	    \,\right)
\\
& = &
    \sigma \cdot \vert\,v_{1} - v_{2}\,\vert \cdot n_{1} n_{2} \cdot
    \left(\,
	    2\cdot\dfrac{AL^{2}}{\vert\,v_{1}-v_{2}\,\vert}
	    \; - \;
	    2\,A \cdot \vert\,v_{1}-v_{2}\,\vert \cdot
	    \left[\,
	        \dfrac{1}{2}\,t^{2}
	        \,\right]_{0}^{L/\vert\,v_{1}-v_{2}\,\vert}
	    \,\right)
\\
& = &
    \sigma \cdot \vert\,v_{1} - v_{2}\,\vert \cdot n_{1} n_{2} \cdot
    \left(\,
	    2\cdot\dfrac{AL^{2}}{\vert\,v_{1}-v_{2}\,\vert}
	    \; - \;
	    A \cdot \vert\,v_{1}-v_{2}\,\vert \cdot
	    \dfrac{L^{2}}{\vert\,v_{1}-v_{2}\,\vert^{2}}
	    \,\right)
\\
& = &
    \sigma \cdot \vert\,v_{1} - v_{2}\,\vert \cdot n_{1} n_{2} \cdot
    \left(\,
        \dfrac{AL^{2}}{\vert\,v_{1}-v_{2}\,\vert}
	    \,\right)
\;\; = \;\;
    \sigma \cdot n_{1} n_{2} \cdot AL^{2}
\\
& = &
    \sigma \cdot \left(n_{1} A L\right) \cdot \left(n_{2} A L\right) \cdot \dfrac{1}{A}
\\
& = &
    \sigma \cdot \dfrac{ N_{1} N_{2} }{A}\,,
\end{eqnarray*}
as required.
\qed

          %%%%% ~~~~~~~~~~~~~~~~~~~~ %%%%%
